\documentclass[conference]{IEEEtran}

\usepackage[T1]{fontenc}
\usepackage[utf8]{inputenc}
\usepackage[spanish,es-nodecimaldot]{babel}
\usepackage{amsmath,amssymb}
\usepackage{graphicx}
\usepackage{booktabs}

\title{Clasificación bayesiana de imágenes por color y diagnóstico de cáncer de mama}

\author{\IEEEauthorblockN{Nombre del estudiante}
\IEEEauthorblockA{Reconocimiento de Patrones}}

\begin{document}
\maketitle
\raggedbottom

\begin{abstract}
Se comparan clasificadores bayesianos en dos problemas: clasificación de píxeles por color y clasificación del conjunto Breast Cancer Wisconsin. Para las imágenes, QDA se ajusta con 1\,440\,000 píxeles de cuatro imágenes de entrenamiento y se analizan cualitativamente tres imágenes de prueba. Para cáncer de mama se entrenan QDA, LDA y Gaussian Naive Bayes (GNB) con una partición estratificada y estandarización. QDA y LDA alcanzan una exactitud de 95.61\% en prueba, frente a 92.98\% de GNB.
\end{abstract}

\begin{IEEEkeywords}
QDA, LDA, Gaussian Naive Bayes, segmentación por color, Breast Cancer Wisconsin.
\end{IEEEkeywords}

\section{Introducción}
Los clasificadores bayesianos asignan una observación a la clase con mayor probabilidad posterior. Su comportamiento depende de los supuestos sobre las distribuciones de las características: QDA estima una matriz de covarianza por clase, LDA supone una covarianza compartida y GNB supone independencia condicional entre características. Esta práctica compara estos métodos en una tarea de clasificación de píxeles y en un problema tabular de diagnóstico, siguiendo la consigna proporcionada [1].

\section{Métodos}
Para una clase $k$, QDA calcula una puntuación discriminante proporcional al logaritmo de la probabilidad posterior:
\begin{equation}
g_k(\mathbf{x})=-\frac{1}{2}(\mathbf{x}-\boldsymbol{\mu}_k)^T\boldsymbol{\Sigma}_k^{-1}(\mathbf{x}-\boldsymbol{\mu}_k)-\frac{1}{2}\log|\boldsymbol{\Sigma}_k|+\log\pi_k,
\end{equation}
donde $\boldsymbol{\mu}_k$, $\boldsymbol{\Sigma}_k$ y $\pi_k$ son, respectivamente, la media, la covarianza y la probabilidad a priori de la clase. La predicción es la clase que maximiza $g_k$.

\section{Parte I: clasificación de imágenes por color}
\subsection{Datos y extracción de características}
Se trabajó con cuatro imágenes de entrenamiento y tres imágenes de prueba, reservadas para evaluar visualmente el clasificador una vez ajustado. Solo las cuatro imágenes de entrenamiento se utilizan para generar las máscaras, construir el DataFrame y estimar los parámetros de QDA.

Las máscaras de plátano, chile y huevo se obtuvieron mediante segmentación en HSV: el verde identifica chile, el amarillo plátano y la baja saturación con alto valor huevo. Los umbrales HSV se usan únicamente para generar etiquetas de entrenamiento; no forman parte del clasificador QDA. Se eliminaron solapamientos entre las máscaras y se aplicaron operaciones morfológicas de apertura y cierre con un elemento estructurante de $3\times3$ píxeles. La máscara de fondo se definió como el complemento de la unión de las tres máscaras de objetos. La Fig.~\ref{fig:mascaras} muestra, para las cuatro imágenes de entrenamiento, la imagen original y las máscaras de plátano, chile, huevo y fondo.

\begin{figure}[!t]
\centering
\includegraphics[width=0.98\columnwidth]{figures/mascaras-entrenamiento-1.png}\\[2pt]
\includegraphics[width=0.98\columnwidth]{figures/mascaras-entrenamiento-2.png}\\[2pt]
\includegraphics[width=0.98\columnwidth]{figures/mascaras-entrenamiento-3.png}\\[2pt]
\includegraphics[width=0.98\columnwidth]{figures/mascaras-entrenamiento-4.png}
\caption{Máscaras generadas para las cuatro imágenes de entrenamiento. Cada panel muestra, de izquierda a derecha, la imagen original, plátano, chile, huevo y fondo.}
\label{fig:mascaras}
\end{figure}

Cada píxel de las máscaras de entrenamiento se representó mediante sus canales rojo, verde y azul y la etiqueta de clase, formando las columnas R, G, B y Clase. El DataFrame resultante tiene dimensiones $1\,440\,000\times4$. El Cuadro~\ref{tab:conteos} resume los píxeles etiquetados por clase, sus probabilidades a priori y el total; el Cuadro~\ref{tab:ejemplos-rgb} presenta tres ejemplos RGB por clase con sus índices originales del DataFrame.

\begin{table}[!t]
\centering
\caption{Píxeles etiquetados y probabilidad a priori por clase.}
\label{tab:conteos}
\small
\begin{tabular}{lrr}
\toprule
Clase & Píxeles & $\pi_k$ \\
\midrule
Fondo & 1\,105\,113 & 0.7674 \\
Plátano & 172\,285 & 0.1196 \\
Huevo & 111\,307 & 0.0773 \\
Chile & 51\,295 & 0.0356 \\
\midrule
Total & 1\,440\,000 & 1.0000 \\
\bottomrule
\end{tabular}
\end{table}

\begin{table}[!t]
\centering
\caption{Ejemplos de píxeles RGB por clase, con índices originales del DataFrame.}
\label{tab:ejemplos-rgb}
\scriptsize
\begin{tabular}{lrrrr}
\toprule
Clase & Índice & R & G & B \\
\midrule
Plátano & 0 & 206 & 138 & 41 \\
Plátano & 1 & 205 & 137 & 38 \\
Plátano & 2 & 212 & 144 & 45 \\
Chile & 43528 & 43 & 54 & 12 \\
Chile & 43529 & 42 & 58 & 11 \\
Chile & 43530 & 44 & 62 & 10 \\
Huevo & 56295 & 251 & 168 & 164 \\
Huevo & 56296 & 243 & 172 & 166 \\
Huevo & 56297 & 247 & 182 & 178 \\
Fondo & 84033 & 222 & 77 & 56 \\
Fondo & 84034 & 220 & 75 & 54 \\
Fondo & 84035 & 219 & 74 & 53 \\
\bottomrule
\end{tabular}
\end{table}

\subsection{Parámetros QDA}
Para cada clase se estimaron la media $\boldsymbol{\mu}_k$, que representa el color RGB promedio; la matriz de covarianza $\boldsymbol{\Sigma}_k$, que describe la dispersión y covariación de los canales RGB; y el prior $\pi_k$, que corresponde a la proporción de píxeles de esa clase en el conjunto de entrenamiento. El Cuadro~\ref{tab:qda-parametros} presenta estas estadísticas, calculadas exclusivamente con las cuatro imágenes de entrenamiento. Las covarianzas se expresan en unidades de intensidad al cuadrado.

\begin{table}[!ht]
\centering
\caption{Estadísticas RGB de QDA por clase.}
\label{tab:qda-parametros}
\scriptsize
\setlength{\tabcolsep}{2pt}
\setlength{\arraycolsep}{1pt}
\renewcommand{\arraystretch}{1.15}
\begin{tabular}{l c c c}
\toprule
Clase & Media $\boldsymbol{\mu}_k$ & Covarianza $\boldsymbol{\Sigma}_k$ & Prior $\pi_k$ \\
\midrule
Chile & $[48.57,\;89.80,\;23.95]$ & $\left[\begin{array}{@{}rrr@{}}522.23&622.29&542.87\\622.29&843.73&681.25\\542.87&681.25&729.62\end{array}\right]$ & 0.0356 \\
\midrule
Fondo & $[218.84,\;91.74,\;61.65]$ & $\left[\begin{array}{@{}rrr@{}}458.69&198.36&197.13\\198.36&168.49&117.79\\197.13&117.79&117.38\end{array}\right]$ & 0.7674 \\
\midrule
Huevo & $[231.54,\;229.88,\;230.85]$ & $\left[\begin{array}{@{}rrr@{}}95.14&116.81&120.67\\116.81&309.39&337.85\\120.67&337.85&384.97\end{array}\right]$ & 0.0773 \\
\midrule
Plátano & $[215.91,\;194.95,\;59.07]$ & $\left[\begin{array}{@{}rrr@{}}323.36&400.95&237.75\\400.95&555.40&333.55\\237.75&333.55&447.51\end{array}\right]$ & 0.1196 \\
\bottomrule
\end{tabular}
\end{table}

QDA emplea una matriz de covarianza diferente para cada clase y asigna cada píxel a la clase con mayor puntuación discriminante según sus características RGB. Por tanto, los umbrales HSV no intervienen en la clasificación posterior.

\subsection{Clasificación y análisis de las imágenes de prueba}
Una vez estimados los parámetros con los datos de entrenamiento, se aplicó QDA a las tres imágenes reservadas. Para cada píxel, sus valores $[R,G,B]$ se introdujeron al clasificador y se obtuvo una clase predicha. Las Figs.~\ref{fig:prueba-1}--\ref{fig:prueba-3} muestran las imágenes de prueba y sus mapas de clasificación.

La Fig.~\ref{fig:prueba-1} muestra la imagen Prueba1 y su mapa QDA.
\begin{figure}[!t]
\centering
\includegraphics[width=0.98\columnwidth]{figures/clasificacion-prueba-1.png}
\caption{Imagen Prueba1 y mapa de clasificación QDA. Negro: fondo; amarillo: plátano; verde: chile; blanco: huevo.}
\label{fig:prueba-1}
\end{figure}

La Fig.~\ref{fig:prueba-2} muestra la imagen Prueba2 y su mapa QDA.
\begin{figure}[!t]
\centering
\includegraphics[width=0.98\columnwidth]{figures/clasificacion-prueba-2.png}
\caption{Imagen Prueba2 y mapa de clasificación QDA. Negro: fondo; amarillo: plátano; verde: chile; blanco: huevo.}
\label{fig:prueba-2}
\end{figure}

La Fig.~\ref{fig:prueba-3} muestra la imagen Prueba3 y su mapa QDA.
\begin{figure}[!t]
\centering
\includegraphics[width=0.98\columnwidth]{figures/clasificacion-prueba-3.png}
\caption{Imagen Prueba3 y mapa de clasificación QDA. Negro: fondo; amarillo: plátano; verde: chile; blanco: huevo.}
\label{fig:prueba-3}
\end{figure}

En los tres mapas se reconocen visualmente las clases principales, aunque aparecen pequeñas confusiones en los bordes. Las sombras, los reflejos y los cambios de iluminación modifican los valores RGB y pueden generar asignaciones entre clases. Además, un objeto que no pertenece a las clases de entrenamiento puede asignarse a la clase cromáticamente más similar. Como no se dispone de máscaras de referencia para las imágenes de prueba, el análisis es cualitativo y no se reporta exactitud por píxel.

\section{Parte II: Breast Cancer Wisconsin}
\subsection{Configuración experimental}
El conjunto completo contiene 569 muestras y 30 características. La clase 0 (maligno) reúne 212 observaciones (37.26\%), y la clase 1 (benigno), 357 (62.74\%). La Fig.~\ref{fig:bc-distribucion} presenta la distribución de las dos clases en el conjunto de datos.

\begin{figure}[!t]
\centering
\includegraphics[width=0.96\columnwidth]{figures/breast-cancer-distribucion-clases.png}
\caption{Distribución de las clases maligno y benigno en Breast Cancer Wisconsin.}
\label{fig:bc-distribucion}
\end{figure}

Se realizó una división estratificada con semilla 42: 455 muestras de entrenamiento (80\%) y 114 de prueba (20\%). En entrenamiento, la clase 0 contiene 170 muestras (37.36\%) y la clase 1, 285 (62.64\%); en prueba, la clase 0 contiene 42 (36.84\%) y la clase 1, 72 (63.16\%). El escalador se ajustó únicamente con entrenamiento y se aplicó a ambos subconjuntos. Se entrenaron QDA con regularización $0.01$, LDA y Gaussian Naive Bayes (GNB).

\subsection{Resultados y análisis}
El Cuadro~\ref{tab:metricas} resume las métricas obtenidas en el conjunto de prueba. La precisión y el recall se presentan por separado para las clases maligno y benigno.

\begin{table}[!t]
\centering
\caption{Resultados en prueba.}
\label{tab:metricas}
\scriptsize
\setlength{\tabcolsep}{2pt}
\begin{tabular}{l c c c c c}
\toprule
Modelo & Accuracy & \shortstack{Precision\\(Maligno)} & \shortstack{Precision\\(Benigno)} & \shortstack{Recall\\(Maligno)} & \shortstack{Recall\\(Benigno)} \\
\midrule
QDA & 95.61\% & 97.44\% & 94.67\% & 90.48\% & 98.61\% \\
LDA & 95.61\% & 97.44\% & 94.67\% & 90.48\% & 98.61\% \\
GNB & 92.98\% & 90.48\% & 94.44\% & 90.48\% & 94.44\% \\
\bottomrule
\end{tabular}
\end{table}

La Fig.~\ref{fig:bc-evaluacion} reúne el gráfico comparativo de accuracy y las matrices de confusión generadas por el cuaderno. En estas matrices, las filas corresponden a la clase real y las columnas a la predicción, en el orden maligno/benigno. QDA y LDA clasificaron correctamente 109 de 114 muestras; GNB clasificó correctamente 106.

\begin{figure}[!t]
\centering
\begin{minipage}[t]{0.47\columnwidth}
\centering
\includegraphics[width=\linewidth]{figures/breast-cancer-comparacion-accuracy.png}
\par\footnotesize (a) Comparación de accuracy.
\end{minipage}\hfill
\begin{minipage}[t]{0.47\columnwidth}
\centering
\includegraphics[width=\linewidth]{figures/breast-cancer-matriz-confusion-qda.png}
\par\footnotesize (b) Matriz de confusión de QDA.
\end{minipage}

\vspace{3pt}

\noindent
\begin{minipage}[t]{0.47\columnwidth}
\centering
\includegraphics[width=\linewidth]{figures/breast-cancer-matriz-confusion-lda.png}
\par\footnotesize (c) Matriz de confusión de LDA.
\end{minipage}\hfill
\begin{minipage}[t]{0.47\columnwidth}
\centering
\includegraphics[width=\linewidth]{figures/breast-cancer-matriz-confusion-gnb.png}
\par\footnotesize (d) Matriz de confusión de GNB.
\end{minipage}
\caption{Evaluación gráfica de QDA, LDA y GNB en el conjunto de prueba.}
\label{fig:bc-evaluacion}
\end{figure}

QDA y LDA presentan el mejor desempeño global y los mismos resultados en esta partición. Ambos identifican 38 de 42 casos malignos y 71 de 72 benignos; sus cuatro falsos negativos corresponden a casos malignos clasificados como benignos. GNB también deja cuatro falsos negativos y aumenta a cuatro los falsos positivos. Por ello, LDA es una elección razonable entre los modelos evaluados: iguala a QDA en las métricas observadas y adopta el supuesto más simple de una covarianza compartida. Sin embargo, el recall maligno de 90.48\% indica que aún se omiten casos; en un contexto clínico, la selección debe ponderar el costo de los falsos negativos, evaluar umbrales y validarse con datos independientes.

\section{Conclusiones}
El QDA basado en color produce mapas visualmente coherentes para los objetos de las tres imágenes, aunque los umbrales HSV y la iluminación generan errores locales. La separación entre entrenamiento y prueba evita que las imágenes de prueba participen en el ajuste; aun así, sin máscaras de referencia solo es posible evaluar los resultados de forma cualitativa. En Breast Cancer Wisconsin, QDA y LDA logran la mayor exactitud (95.61\%) y superan a GNB (92.98\%). Dado el empate observado, LDA ofrece un modelo más simple, pero la sensibilidad a casos malignos y la validación en otras particiones deben guiar cualquier uso posterior.

\begin{thebibliography}{1}
\bibitem{practica}
\emph{Práctica 3: Clasificadores Bayesianos}. Material de la asignatura Reconocimiento de Patrones, documento y cuadernos de trabajo proporcionados para la práctica.
\end{thebibliography}

\end{document}
